These notes are an introduction to probabilistic thinking: the mathematical art of reasoning under uncertainty. Probability is not presented here as a collection of ready-made formulas, but as a language for describing experiments, incomplete information, variability, randomness, and data.

The purpose of the course is to teach how to build probabilistic models carefully. Before one computes a probability, expectation, variance, confidence interval, or approximation, one should understand what is random, what is being observed, what assumptions are being made, and what mathematical object represents the situation. In this sense, probability is not only a computational tool, but also a discipline of precise modelling.

The course is intended for readers who want to understand why probabilistic methods work, not only how to apply them. The emphasis is therefore placed on interpretation, model construction, and the connection between formal definitions and concrete situations. Calculations appear throughout the notes, but they are treated as consequences of a well-defined model rather than as isolated procedures.

\section{How to Use These Notes}

These notes are meant to be read before problem classes. A problem class is not the place where the basic language of the topic is introduced for the first time; it is the place where that language is tested, used, questioned, and sharpened. Students are therefore expected to read the relevant chapter in advance, think through the examples and explanations, and come prepared to discuss the ideas that appear in the text.

Reading these notes should be an active process. The goal is not to recognize formulas, but to understand the modelling decisions behind them. While reading, students should repeatedly ask: What is the experiment? What counts as an elementary outcome? What information is given? What is random? What assumptions are being made? What event, random variable, distribution, or statistic is being constructed? Why is this model appropriate, and what would change if the assumptions changed?

Problem lists, computations, and artificial intelligence tools may be used as support, and students are encouraged to use them responsibly when they help with practice, checking ideas, generating additional examples, or clarifying difficult points. However, such tools do not replace the central task: understanding the notes. A correct numerical answer is not sufficient if the student cannot explain the model, justify the method, interpret the result, and answer questions about the reasoning used in the chapter.

The expected preparation for classes is therefore conceptual as well as computational. After reading a chapter, a student should be able to summarize its main ideas, explain the definitions in their own words, reconstruct the key examples, identify the assumptions in a probabilistic model, and answer questions that test understanding of the text. The notes are deliberately written with full explanations and detailed argumentation so that this preparation is possible.

\section{Structure of the Course}

These lecture notes are organized as a gradual construction of probabilistic thinking. The aim is not to begin with formulas and then search for examples, but to develop the mathematical language needed to describe uncertainty in a precise way. We start from the problem of counting and representing possible outcomes, then move toward events, probability, random variables, probability laws, empirical data, sampling, limit theorems, and statistical inference.

The course begins with combinatorics because many probabilistic models depend on the ability to count possible outcomes correctly. At this stage probability itself is not yet needed. What is needed is the more elementary skill of recognizing what kind of object an outcome is: a sequence, a selection, an arrangement, a choice with or without repetition, or an object whose elements may or may not be distinguishable. This first step builds the counting language that later allows us to describe sample spaces and favorable outcomes with precision.

The second stage has a special conceptual role. It is devoted to discovering the formal language of probability theory. The purpose is to show how a mathematical theory can be built from careful thinking, analysis, and abstraction. We examine how statements about a random situation select sets of outcomes, how these sets become events, how logical operations correspond to set operations, and why a numerical theory of probability should be placed on such a structure. This part is not primarily computational; it is meant to teach how probabilistic language is formed.

After this conceptual preparation, the notes turn to sample spaces and elementary probability. Several canonical examples are developed from the beginning: one first understands the experiment, decides what should count as an elementary outcome, constructs the sample space, defines events, and only then assigns probabilities. Different representations of the same structure, such as lists, tables, trees, and symbolic notation, are used to show that a sample space is an abstract mathematical model rather than a collection of observed data.

The next step is conditional probability. Here the central idea is that information changes the space in which we reason. Conditional probability is introduced by asking what it means to restrict attention to the cases in which some event has already occurred. From this point of view, the formula for conditional probability becomes natural, and it leads to independence, dependence, partitions, the law of total probability, and Bayes' theorem. Bayes' theorem is treated as a central tool for reversing conditional reasoning and for understanding situations where intuition is easily misled.

The course then introduces random variables through the discrete case. A random variable appears as a function that extracts a numerical quantity from an elementary outcome: the number of successes, the number of errors, the number of arrivals, a payoff, or another measurable feature of the experiment. This leads to the support of a discrete random variable, its probability mass function, its cumulative distribution function, expectation, variance, moments, and interpretation. Classical discrete laws are presented as models arising from specific mechanisms, not as isolated formulas.

The continuous case is then developed as a parallel but conceptually more demanding part of the theory. For a continuous random variable, probability is no longer concentrated at individual points, so probabilities of intervals become the fundamental objects. This leads to density functions, cumulative distribution functions, integrals, continuous expectations, variances, moments, quantiles, and the interpretation of probability as area under a density. Important continuous models, including the uniform, exponential, normal, gamma, and beta distributions, are introduced with attention to their meaning and use.

Once theoretical probability laws have been introduced, the course moves to empirical distributions and descriptive statistics. The same ideas now appear in data: frequency tables, histograms, empirical cumulative distribution functions, means, medians, variances, standard deviations, quartiles, interquartile ranges, and outliers. This stage is meant to clarify the difference between an idealized probability law and a finite dataset, and also to show how empirical summaries mirror the theoretical quantities defined earlier.

The following stage adds another level of randomness by studying sampling and distributions of statistics. A sample mean, sample proportion, sample variance, median, or other statistic is not merely a number; it is itself a random variable because it depends on the random sample. This perspective prepares the ground for understanding sampling distributions and for seeing why repeated sampling creates a new layer of probabilistic structure.

The final stage of the present course outline is devoted to limit theorems and statistical inference. Once empirical summaries and sampling distributions are understood, we can ask what regularities appear as the sample size grows or as experiments are repeated many times. This leads to the law of large numbers, the central limit theorem, normal approximations, standard errors, confidence intervals, and the basic logic of statistical inference. The goal is to show how probability theory makes inference from data possible: it allows us to estimate unknown quantities, measure uncertainty, and justify conclusions mathematically.
